Es ist

\mavergleichskettedisp
{\vergleichskette
{ h(x,y,z)
}
{ =} { x^2+y^2 -2 \sqrt{x^2+y^2}R+z^2 + R^2
}
{ } {
}
{ } {
}
{ } {
}
}
{}{}{}
und

\mavergleichskette
{\vergleichskette
{ T
}
{ = }{ h^{-1} (r^2)
}
{  }{
}
{    }{
}
{   }{
}
}
{}{}{.}
Der Gradient von $h$ ist

\mavergleichskettedisp
{\vergleichskette
{
\operatorname{Grad} \,  h
}
{ =} { 2 \begin{pmatrix}   x   -  xR \left(  x^2+y^2  \right)^{- { \frac{   1  }{  2 } } }  \\ y   -  y R \left(  x^2+y^2  \right)^{- { \frac{   1  }{  2 } } } \\ z   \end{pmatrix}
}
{ } {
}
{ } {
}
{ } {
}
}
{}{}{.}
Das
\definitionsverweis {Einheitsnormalenfeld}{}{}
ist daher

\mavergleichskettedisp
{\vergleichskette
{ N(x,y,z)
}
{ =} { { \frac{   1  }{  \sqrt{ \left(   x   -  xR \left(  x^2+y^2  \right)^{- { \frac{   1  }{  2 } } }  \right)^2 + \left(  y   -  y R \left(  x^2+y^2  \right)^{- { \frac{   1  }{  2 } } }  \right)^2   +z^2   }  } } \begin{pmatrix}   x   -  xR \left(  x^2+y^2  \right)^{- { \frac{   1  }{  2 } } }  \\ y   -  y R \left(  x^2+y^2  \right)^{- { \frac{   1  }{  2 } } } \\ z   \end{pmatrix}
}
{ } {
}
{ } {
}
{ } {
}
}
{}{}{.}

 [[Kategorie:Latexseite]]
